\documentclass[11pt]{amsart}

\usepackage[T1]{fontenc}
\usepackage{lmodern}
\usepackage{microtype}
\usepackage{mathtools,amssymb,amsthm}
\usepackage[hidelinks]{hyperref}

\hypersetup{
  pdftitle={A characteristic-free integral witness for Froberg's series for six
            to twenty-one septics in four variables}
}

\newtheorem{theorem}{Theorem}
\newtheorem{lemma}[theorem]{Lemma}
\newtheorem{proposition}[theorem]{Proposition}
\newtheorem{corollary}[theorem]{Corollary}
\theoremstyle{remark}
\newtheorem*{remark}{Remark}

\DeclareMathOperator{\rank}{rank}
\newcommand{\NN}{\mathbb N}
\newcommand{\ZZ}{\mathbb Z}
\newcommand{\QQ}{\mathbb Q}
\newcommand{\FF}{\mathbb F}

\title[A characteristic-free witness for six to twenty-one septics]
{A characteristic-free integral witness for Fr\"oberg's series for six to
twenty-one septics in four variables}

\date{}

\subjclass[2020]{13D40, 13A02, 13P10, 15A15}
\keywords{Fr\"oberg's conjecture, Hilbert series, generic forms, positive
characteristic, determinantal divisor, coprime minors}

\begin{document}

\begin{abstract}
Fr\"oberg's Conjecture~1, as stated in the Fr\"oberg--Lundqvist survey, predicts
the Hilbert series of a quotient of a polynomial ring by generic forms; the
positive-characteristic cases the survey records are $n\le 2$ and $n=3$ over an
infinite field. We exhibit, for every $r$ with $6\le r\le 21$, an explicit
$r$-tuple of septics in $\ZZ[x_1,x_2,x_3,x_4]$ with all coefficients $0$ or $1$
whose reduction modulo the characteristic of a field $k$ generates an ideal of
$k[x_1,x_2,x_3,x_4]$ with the conjectured quotient series, for \emph{every}
field $k$. This is a statement about one explicit tuple and not about generic
forms; specialisation carries it to the generic forms over the purely
transcendental coefficient field $k(c_{i,a})$, so Conjecture~1 for
$(n;d_1,\dots,d_r)=(4;7,\dots,7)$ and $6\le r\le21$ holds over such a field in
every characteristic, and nothing is claimed over a finite field, where there
are no generic forms in that sense. The certificate is a coprime pair: for
each of ten endpoint cells $(r,j)$ we exhibit two maximal square minors of one
integer matrix whose determinants are coprime, so no prime divides both, so the
rank is maximal modulo every prime at once; ten cells then spread over the whole
interval by a characteristic-free propagation lemma, proved here. The family,
the endpoint table and the propagation lemma are Wang--Zhang's, as is the
characteristic-zero case; inside $6\le r\le21$ the characteristic-$2$ slice is
already an immediate consequence of their printed rank table over $\FF_2$
together with their propagation lemma, which carries no field hypothesis. What
is added here is the coprime-pair certificate and, with it, the remaining
primes.
\end{abstract}

\maketitle

\section{The conjecture}

Let $k$ be a field, $S=k[x_1,\dots,x_n]$ with the standard grading, and let
$f_1,\dots,f_r$ be forms of degrees $d_1,\dots,d_r$. Write
$R=S/(f_1,\dots,f_r)$. Fr\"oberg's conjecture predicts the Hilbert series of $R$
when the $f_i$ are \emph{generic}, meaning that their coefficients are
algebraically independent over the prime field of $k$. In the numbering of the
Fr\"oberg--Lundqvist survey \cite[Conjecture~1]{FL} the statement reads,
verbatim:

\begin{quote}
$R(z)=\Bigl(\prod_{i=1}^{r}(1-z^{d_i})/(1-z)^n\Bigr)_+$, \emph{where}
$\bigl(\sum_{i\ge0}a_iz^i\bigr)_+=\sum_{i\ge0}b_iz^i$, \emph{where} $b_i=a_i$
\emph{if} $a_j\ge0$ \emph{for all} $j\le i$ \emph{and} $b_i=0$
\emph{otherwise.}
\end{quote}

\noindent
Conjecture~1 sits on p.~416 of \cite{FL}. The survey assumes $k=\mathbb C$
unless otherwise stated (p.~415); its Section~7, \emph{Positive characteristic}
(p.~426), is where it is otherwise stated, and the positive-characteristic cases
recorded there are $n\le2$ together with $n=3$ over an infinite field, followed
by the sentence ``We believe that Conjecture~1 is true for any field.'' So the
conjecture is posed over an arbitrary field by its own authors. For $n=4$ and
$d=7$ in characteristic $p$, $r\le4$ is the monomial complete intersection
$(x_1^7,\dots,x_4^7)$.

\section{The result}

Throughout, $n=4$, $d=7$, $S=k[x_1,x_2,x_3,x_4]$, and for forms
$F_1,\dots,F_r\in S_7$ we write $I(r)=(F_1,\dots,F_r)$ and
\[
  \mu_{r,j}\colon (S_{j-7})^{r}\longrightarrow S_j,
  \qquad (h_1,\dots,h_r)\longmapsto \sum_{i=1}^{r}h_iF_i ,
\]
whose image is $I(r)_j$, so that
\begin{equation}\label{eq:hf}
  \dim_k (S/I(r))_j=\binom{j+3}{3}-\rank \mu_{r,j}.
\end{equation}
For the equal-degree case at hand put
\[
  a_j(r)=[z^j]\;\frac{(1-z^{7})^{r}}{(1-z)^{4}}
        =\sum_{i\ge0}(-1)^i\binom{r}{i}\binom{j-7i+3}{3},
\]
and let $q(r)$ be the largest $j$ with $a_0(r),\dots,a_j(r)$ all positive. The
truncation in Conjecture~1 then says: the conjectured series has coefficient
$a_j(r)$ for $j\le q(r)$ and $0$ for $j>q(r)$. (If some $a_j(r)$ vanishes, the
two readings of ``$a_j\ge0$'' in the conjecture's truncation rule agree, because
the coefficient printed is $0$ either way; for $6\le r\le21$ no coefficient
$a_j(r)$ with $j\le q(r)$ vanishes, and this is re-derived in the verification.)

\begin{theorem}\label{thm:main}
Let $k$ be \emph{any} field and $S=k[x_1,x_2,x_3,x_4]$. Let
$G_1,\dots,G_{21}\in\ZZ[x_1,x_2,x_3,x_4]_7$ be the $21$ septics printed in
Appendix~\ref{app:family}, each a sum of distinct degree-$7$ monomials with
coefficient $1$, and read them in $S$ by reducing their coefficients modulo the
characteristic of $k$. Then for every $r$ with $6\le r\le 21$ the quotient
$S/(G_1,\dots,G_r)$ has Hilbert series
\[
  \sum_{j=0}^{q(r)}a_j(r)\,z^{j}
  \;=\;\Bigl(\tfrac{(1-z^{7})^{r}}{(1-z)^{4}}\Bigr)_+ ,
\]
the series predicted by Conjecture~1 for $(n;d_1,\dots,d_r)=(4;7,\dots,7)$.
\end{theorem}

Theorem~\ref{thm:main} is a statement about one explicit integral tuple; no
genericity enters it. Genericity enters only through the specialisation argument
below, which is the mechanism behind the criterion that \cite[p.~416]{FL} states
for this purpose: ``if one finds an algebra with the conjectured series for some
$(n,r,d_1,\dots,d_r)$, then the conjecture is proved for these values''.

\begin{proposition}\label{prop:generic}
Let $k$ be a field, $K=k(c_{i,a})$ the purely transcendental extension of $k$
generated by indeterminates $c_{i,a}$ indexed by $1\le i\le r$ and the
monomials $x^a$ of degree $7$, and $F_i=\sum_a c_{i,a}x^a\in K[x_1,\dots,x_4]_7$
the generic forms. If some specialisation of the $c_{i,a}$ to elements of $k$
makes $\rank\mu_{r,j}$ equal to $\min\bigl(\dim S_j,\;r\dim S_{j-7}\bigr)$ for
every $j$, then the same holds for the generic $F_i$ over $K$, and hence
$K[x]/(F_1,\dots,F_r)$ has the series of Theorem~\ref{thm:main}.
\end{proposition}

\begin{proof}
The entries of the matrix of $\mu_{r,j}$ are $\ZZ$-linear in the $c_{i,a}$. Let
$s=\min(\dim S_j,\;r\dim S_{j-7})$. By hypothesis some $s\times s$ minor is
nonzero after the specialisation; that minor is a polynomial in the $c_{i,a}$
with coefficients in the prime field, and it does not vanish at a point of
$k^N$, hence it is not the zero polynomial, hence it is nonzero in $K$. So the
generic rank is at least $s$, and it is at most $s$ trivially. Now
\eqref{eq:hf} computes the generic Hilbert function.
\end{proof}

Theorem~\ref{thm:main} supplies the hypothesis of
Proposition~\ref{prop:generic} over a field of any characteristic: for the
specialisation sending $c_{i,a}$ to the coefficients of $G_i$, its dimension
formula read through \eqref{eq:hf} says exactly that $\rank\mu_{r,j}$ is the
smaller of $\dim S_j$ and $r\dim S_{j-7}$ in every degree $j$. So for every
field $k$ Conjecture~1 holds for $n=4$, $d_1=\dots=d_r=7$ and $6\le r\le21$ over
the coefficient field $K=k(c_{i,a})$, whose characteristic is that of $k$ and
hence arbitrary. Over a \emph{finite} field there are no generic forms in the
above sense, and we claim nothing about Conjecture~1 there; what
Theorem~\ref{thm:main} gives over $\FF_p$ is the explicit tuple itself.

\section{The witness family}\label{sec:family}

Order the $\binom{10}{3}=120$ monomials of degree $7$ in $x_1,x_2,x_3,x_4$ by
\emph{decreasing lexicographic order of their exponent vectors}, so that the
list begins $[7,0,0,0]$, $[6,1,0,0]$, $[6,0,1,0]$, $[6,0,0,1]$, $[5,2,0,0]$,
\dots\ and ends $[0,0,1,6]$, $[0,0,0,7]$; call this ordered list
$\mathrm{MON}_7$ and index it from $0$. Each form of the family is a sum of
distinct monomials from $\mathrm{MON}_7$, all coefficients $1$, and is written as
a $30$-digit hexadecimal \emph{mask}: the monomial with index $i$ occurs in the
form if and only if bit $i$ --- the coefficient of $2^i$ --- of the mask is $1$.
The family has $120$ members; Appendix~\ref{app:family} prints the masks of the
$21$ used in Theorem~\ref{thm:main}, whose support sizes run from $49$ to $78$,
and all $120$ masks are listed in the accompanying program \texttt{verify.py}.

To make the convention checkable by hand, here is $G_1$ in full. Its mask is
\texttt{49ed660e53a564e3fb9f22d51e14f7} and it is the sum of $x^a$ over these $65$ exponent
vectors $a$:

\begin{center}\footnotesize\begin{verbatim}
[7,0,0,0] [6,1,0,0] [6,0,1,0] [5,2,0,0] [5,1,1,0] [5,1,0,1]
[5,0,2,0] [4,3,0,0] [4,2,0,1] [4,0,2,1] [4,0,1,2] [4,0,0,3]
[3,4,0,0] [3,2,1,1] [3,1,3,0] [3,1,1,2] [3,0,4,0] [3,0,3,1]
[3,0,1,3] [2,4,0,1] [2,3,0,2] [2,2,3,0] [2,2,2,1] [2,2,1,2]
[2,2,0,3] [2,1,2,2] [2,1,1,3] [2,1,0,4] [2,0,4,1] [2,0,3,2]
[2,0,2,3] [2,0,1,4] [2,0,0,5] [1,6,0,0] [1,5,1,0] [1,4,0,2]
[1,3,3,0] [1,3,2,1] [1,2,4,0] [1,2,1,3] [1,2,0,4] [1,1,4,1]
[1,1,2,3] [1,0,6,0] [1,0,4,2] [1,0,3,3] [1,0,2,4] [0,7,0,0]
[0,6,0,1] [0,5,0,2] [0,4,3,0] [0,4,2,1] [0,3,1,3] [0,3,0,4]
[0,2,3,2] [0,2,2,3] [0,2,0,5] [0,1,5,1] [0,1,4,2] [0,1,2,4]
[0,1,1,5] [0,1,0,6] [0,0,7,0] [0,0,4,3] [0,0,1,6]
\end{verbatim}\end{center}

The witness for a given $r$ is the \emph{prefix} $G_1,\dots,G_r$. Writing each
form as the sorted list of its exponent vectors and the family of $120$ as the
list of those lists, the canonical JSON encoding
\texttt{[[[7,0,0,0],[6,1,0,0],\dots]],\dots]} (no whitespace) has
\[
  \text{\texttt{sha256}} =
  \text{\footnotesize\texttt{f93339e7b30deb19b81fe74fa3b79a7859cf8da9a4c6fadbd740f98c975c8e20}}.
\]
This is the digest of the \texttt{forms} array of the ancillary certificate
published with \cite{WZ}: the family is theirs.

\section{Ten cells, and why they suffice}

\begin{lemma}[endpoint propagation; \cite{WZ}]\label{lem:endpoint}
Let $k$ be any field, $S=k[x_1,\dots,x_n]$, and $G_1,\dots,G_b\in S_d$. Let
$1\le a\le b$ and $q\ge d$, and suppose
\begin{enumerate}
\item[(i)] $\mu_{b,q}$ is injective, and
\item[(ii)] $\mu_{a,q+1}$ is surjective.
\end{enumerate}
Then for every $r$ with $a\le r\le b$: $\mu_{r,j}$ is injective for
$d\le j\le q$, and $(S/I(r))_j=0$ for every $j\ge q+1$. Consequently
\[
  \dim_k (S/I(r))_j=
  \begin{cases}
   \binom{j+n-1}{n-1}, & 0\le j<d,\\[2pt]
   \binom{j+n-1}{n-1}-r\binom{j-d+n-1}{n-1}, & d\le j\le q,\\[2pt]
   0,& j>q.
  \end{cases}
\]
\end{lemma}

\begin{proof}
(a) For $r\le b$, $\mu_{r,q}$ is the restriction of $\mu_{b,q}$ to the subspace
$\{h_{r+1}=\dots=h_b=0\}$, so it is injective by (i).
(b) Let $d\le j\le q$ and let $h=(h_1,\dots,h_r)$ satisfy $\mu_{r,j}(h)=0$ with
$h\ne0$. Choose any monomial $m$ of degree $q-j$. Then $mh\ne0$, because $S$ is
a domain, and $\mu_{r,q}(mh)=m\,\mu_{r,j}(h)=0$, contradicting (a). So
$\mu_{r,j}$ is injective.
(c) For $r\ge a$ the image of $\mu_{r,q+1}$ contains the image of
$\mu_{a,q+1}$, which is all of $S_{q+1}$ by (ii); hence $I(r)_{q+1}=S_{q+1}$.
(d) For $j\ge q+1$, $S_j=S_{j-q-1}\cdot S_{q+1}=S_{j-q-1}\cdot I(r)_{q+1}
\subseteq I(r)_j$, so $(S/I(r))_j=0$.
The displayed dimensions are then \eqref{eq:hf} with
$\rank\mu_{r,j}=r\dim S_{j-d}$ in the injective range. No hypothesis on $k$ is
used beyond $S$ being a domain, which holds over every field; in particular the
lemma transports to characteristic $p$ unchanged.
\end{proof}

For $6\le r\le21$ the value $q(r)$ is constant on five blocks:
\[
\begin{array}{c|ccccc}
 \text{block of } r & \{6\} & \{7,8\} & \{9,10\} & \{11,\dots,14\}
   & \{15,\dots,21\}\\\hline
 q(r) & 13 & 12 & 11 & 10 & 9
\end{array}
\]
so Lemma~\ref{lem:endpoint} needs one injective endpoint $(b,q)$ and one
surjective endpoint $(a,q+1)$ per block: ten cells in all, listed in Table~1.
The blocks cover $\{6,\dots,21\}$ with no gap and no overlap, and the arithmetic
behind the table --- including that no second Koszul term
$\binom{r}{2}\binom{j-14+3}{3}$ enters $a_j(r)$ in the ranges used, the one
exception being the surjective endpoint $(6,14)$, where
$a_{14}(6)=680-720+15=-25<0$ so that the truncation still begins at $q+1$ --- is
re-derived in the verification.

\begin{table}[ht]
\caption{The ten endpoint cells. In each row the matrix is that of $\mu_{r,j}$
in the monomial bases, of size $\dim S_j\times r\dim S_{j-7}$;
$s=\min$ of the two, so ``injective'' means $\rank=s=\#\text{columns}$ and
``surjective'' means $\rank=s=\#\text{rows}$. The last two columns give the bit
lengths of the two exhibited minor determinants and their gcd.}
\begin{center}
\begin{tabular}{ccccccc}
$r$ & $j$ & matrix & $s$ & endpoint & $\log_2$ of the two $|\det|$ & gcd\\\hline
$6$ & $13$ & $560\times 504$ & $504$ & injective & $800$, $789$ & $1$ \\
$6$ & $14$ & $680\times 720$ & $680$ & surjective & $1070$, $1073$ & $1$ \\
$8$ & $12$ & $455\times 448$ & $448$ & injective & $745$, $737$ & $1$ \\
$7$ & $13$ & $560\times 588$ & $560$ & surjective & $905$, $912$ & $1$ \\
$10$ & $11$ & $364\times 350$ & $350$ & injective & $592$, $570$ & $1$ \\
$9$ & $12$ & $455\times 504$ & $455$ & surjective & $756$, $752$ & $1$ \\
$14$ & $10$ & $286\times 280$ & $280$ & injective & $475$, $467$ & $1$ \\
$11$ & $11$ & $364\times 385$ & $364$ & surjective & $611$, $613$ & $1$ \\
$21$ & $9$ & $220\times 210$ & $210$ & injective & $365$, $362$ & $1$ \\
$15$ & $10$ & $286\times 300$ & $286$ & surjective & $484$, $488$ & $1$ \\
\end{tabular}
\end{center}
\end{table}

\begin{lemma}\label{lem:puniform}
Let $M$ be an $R\times C$ integer matrix and $s=\min(R,C)$. Suppose two
$s\times s$ submatrices $N_1,N_2$ of $M$ satisfy
$\gcd(\det N_1,\det N_2)=1$. Then for every field $k$ the image of $M$ in
$k^{R\times C}$ has rank exactly $s$.
\end{lemma}

\begin{proof}
If $\operatorname{char}k=p>0$ then $p$ cannot divide both determinants, so one
of them is a unit in $\FF_p\subseteq k$; if $\operatorname{char}k=0$ then
$\gcd(\det N_1,\det N_2)=1$ forces $\det N_1\ne0$, and a nonzero integer is
nonzero in $k$. Either way some $s\times s$ minor of the image of $M$ is
nonzero, so its rank is at least $s$, and it is at most $s$ by size.
\end{proof}

Lemma~\ref{lem:puniform} is the whole of the new input. A single maximal minor
that is nonzero modulo one prime --- for instance nonzero modulo $2$, which is
what the rank tables of \cite{WZ} record --- certifies that the integer
determinant is nonzero, hence maximal rank over $\QQ$ and over every field of
characteristic zero, and \emph{nothing} at any other prime. Two maximal minors
of the same integer matrix with coprime determinants certify maximal rank at
every prime at once.

\begin{proof}[Proof of Theorem~\ref{thm:main}]
For each of the ten cells $(r,j)$ of Table~1 the accompanying program
\texttt{verify.py} holds two index sets, each naming an $s\times s$ submatrix of
the matrix of $\mu_{r,j}$ built from the prefix $G_1,\dots,G_r$ over $\ZZ$,
together with the two integer determinants, whose bit lengths and gcd are the
last two columns of Table~1. Rows there are indexed from $0$ by $\mathrm{MON}_j$
in decreasing lexicographic order, and columns from $0$ so that column
$i\cdot\dim S_{j-7}+m$ carries the $m$-th monomial of $\mathrm{MON}_{j-7}$
multiplying $G_{i+1}$, with $0\le i<r$; at an injective endpoint the minor uses
every column and the rows outside the listed set, at a surjective endpoint every
row and the columns outside it. The two determinants are coprime in every one of
the ten cells. By Lemma~\ref{lem:puniform} the rank of $\mu_{r,j}$ is $s$ over every
field $k$, which is injectivity at the five cells with $s=\#\text{columns}$ and
surjectivity at the five cells with $s=\#\text{rows}$. Feeding the two endpoints
of each block into Lemma~\ref{lem:endpoint} gives, for every $r$ in that block
and every field $k$, the dimensions displayed there; and those dimensions are
$a_j(r)$ for $j\le q(r)$ and $0$ beyond, which is the truncated series.
\end{proof}

\section{Scope, and what is taken from elsewhere}

\noindent\textbf{The interval.} For $r=21$ one has $a_9(21)=10>0$ and
$a_{10}(21)=-134$, so the predicted series has support through degree~$9$ and
vanishes from degree~$10$ on; establishing the conjecture at $r=21$ accordingly
requires surjectivity in degree~$10$, which is what the surjective endpoint
$(15,10)$ of Table~1 supplies through Lemma~\ref{lem:endpoint}. For $r=22$ one
has $a_9(22)=0$, so the predicted series
stops at degree~$8$. Nothing is claimed here for $r\ge22$, and nothing for
$r=5$; at the lower end $r\le4$ is the monomial complete intersection.

\medskip\noindent\textbf{What is taken from \cite{WZ}, and what is added.} The
family of Appendix~\ref{app:family} is theirs, the block table is theirs, and
Lemma~\ref{lem:endpoint} is theirs; the second is re-derived here from the series
arithmetic, and the third is proved above rather than cited so that the argument
is self-contained. Their Theorem~1.1 settles the characteristic-zero case, and
their Table~3 records maximal rank over $\FF_2$ for the ten cells of Table~1.
Since their propagation lemma is stated with no hypothesis on the field, the
characteristic-$2$ slice of Theorem~\ref{thm:main} for $6\le r\le21$ is an
immediate consequence of that table together with that lemma; they do not assert
it, their Section~5 claiming characteristic zero only. What is added here is
Lemma~\ref{lem:puniform} and the coprime pairs, which cover the remaining
primes, and the fact that the witness is a single explicit tuple over $\FF_p$
rather than an open set of tuples. The integer determinants were computed
independently of \cite{WZ},
which reports ranks modulo $2$ and no integer determinants.

\section{Verification}

The folder that accompanies this note contains one program, \texttt{verify.py}
(Python~3.9 or later, standard library only, no third-party package and no
external data file), and a recorded run of it, \texttt{verify.output.txt}. The
program holds the $120$ masks, the ten cells of Table~1 and the twenty index sets
and twenty integer determinants, and re-derives every quantity claimed here: it
decodes the masks and recomputes the \texttt{sha256} of \S\ref{sec:family};
re-derives $q(r)$ for $r=6,\dots,21$, the five blocks, and the claim that the ten
cells of Table~1 are exactly the ten endpoint cells; rebuilds each of the ten
integer matrices from the decoded forms, checks its shape against $\dim S_j$ and
$r\dim S_{j-7}$, extracts the twenty named minors and recomputes their
determinants from scratch; takes the ten gcds; and, independently of the gcd
argument, recomputes the rank of each of the ten matrices over $\FF_2$, $\FF_3$
and $\FF_5$. The arithmetic is exact: modular Gaussian elimination over $62$-bit
prime fields and Chinese remaindering against the integer Hadamard bound
$\det(N)^2\le\prod_i(\text{$i$-th row sum of }N)$, valid because every entry is
$0$ or $1$, so that the modulus used exceeds $2|\det N|$ and the lift is the true
integer; then one Euclid gcd per cell. No floating point enters any decision. The
recorded run reports
\begin{center}
\texttt{VERDICT: ALL 222 CHECKS PASS}
\end{center}
and the program exits $0$ if and only if every check passes.

Two limits are recorded in the program's own output. It reads no file and no
network, so it does not fetch the ancillary file of \cite{WZ}: the
\texttt{sha256} of \S\ref{sec:family} is recomputed from the masks it holds, and
its equality with the published digest is a claim a referee must check against
arXiv. And it checks nothing about $r\le5$, $r\ge22$, or any $(n,d)$ other than
$(4,7)$.

\appendix

\section{The twenty-one forms}\label{app:family}

The $21$ forms used in Theorem~\ref{thm:main}, as $30$-digit hexadecimal masks
over the ordered basis $\mathrm{MON}_7$ of \S\ref{sec:family}. Bit $i$ of the
mask is $1$ if and only if the $i$-th monomial of $\mathrm{MON}_7$ occurs in the
form.

{\scriptsize\begin{verbatim}
G1    49ed660e53a564e3fb9f22d51e14f7
G2    2728e9577f9976d66df364ad739e9b
G3    e4642c6e3d0c00c837aad525608035
G4    a697391e0adb52a5bbbdb4606db8a6
G5    476783c09aedb8c07d24678ea4bbdd
G6    b29ddfdd727b6b736069bafd78deff
G7    53c2fe37a5bc1a5e7c15ddfefcc5a3
G8    8bc155a7fb2fdaca14f9b5219eeaf4
G9    6703cbffc1ee6058f28ea9114ed06c
G10   4c1f42a4ceb4857b5f07e66cedd8ef
G11   4cc658bc60373a0e6d46cc4493e1c5
G12   f380fa591a18fe94d6e091acd9979a
G13   d654523b92b9bb4034ce8b4d293176
G14   44e3a772aea91b2853a289bd834cc9
G15   2b2671f51fac829ce9e991ce47bc38
G16   b53b6f0b4b896b8cbf7f55fb902070
G17   4495e190e5533494c2071c137aa59d
G18   73837221ac7cb93029c29f2e30e086
G19   26ed68177571506497c38afb7b43fd
G20   42be9cc2a6dbe2f55faf1411447aa5
G21   9fe550885e47ed36507d97205539e0
\end{verbatim}}

\begin{thebibliography}{9}

\bibitem{FL}
R.~Fr\"oberg and S.~Lundqvist,
\emph{Questions and conjectures on extremal Hilbert series},
Rev.\ Un.\ Mat.\ Argentina \textbf{59} (2018), no.~2, 415--429.
\href{https://doi.org/10.33044/revuma.v59n2a10}{doi:10.33044/revuma.v59n2a10}.
All quotations above are from this text.

\bibitem{WZ}
Y.~Wang and L.~Zhang,
\emph{Fr\"oberg's conjecture for quintics and septics in four variables},
arXiv:2608.24797.
The origin of the family of Appendix~\ref{app:family}, of the block table, and
of Lemma~\ref{lem:endpoint}. A preprint, not refereed.

\end{thebibliography}

\end{document}
